\documentclass[11pt]{article}
\usepackage[margin=1in]{geometry}
\usepackage{amsmath,amssymb}
\usepackage[T1]{fontenc}
\usepackage[utf8]{inputenc}
\usepackage{lmodern}
\usepackage[unicode]{hyperref}

\title{AMReX FFT Poisson Benchmark Notes}
\author{Diego Juarez}
\date{April 2026}

\begin{document}
\maketitle

\section*{Purpose}
This note explains the physics problem and the computational workflow behind the AMReX benchmark located in
\texttt{codes/amrex/fft\_poisson\_benchmark}. The goal of that mini-app is not to show off adaptive mesh refinement yet. The goal is narrower and cleaner: reproduce the same manufactured periodic Poisson problem already solved by the repository's standalone C++ FFT code, but now through AMReX's \texttt{FFT::Poisson} interface so that the result can be compared in a controlled way inside the notebook
\texttt{notebooks/AMReX Comparison - 3D FFT Poisson Benchmark.ipynb}.

\section{Physics Problem}
The electrostatic potential $\Phi(\mathbf r)$ satisfies Poisson's equation
\[
\nabla^2 \Phi = -\frac{\rho}{\varepsilon_0}.
\]
For this benchmark we do not start from a difficult physical charge distribution and hope the solver is right. Instead, we use a manufactured solution: choose an exact potential first, then derive the source term from it. That way the numerical answer can be checked point-by-point.

The exact periodic potential is chosen as
\[
\Phi_{\text{exact}}(x,y,z)
=
\cos\!\left(\frac{2\pi x}{L_x}\right)
\cos\!\left(\frac{2\pi y}{L_y}\right)
\cos\!\left(\frac{2\pi z}{L_z}\right).
\]
Applying the Laplacian gives
\[
\nabla^2 \Phi_{\text{exact}}
=
-\left[
\left(\frac{2\pi}{L_x}\right)^2
+
\left(\frac{2\pi}{L_y}\right)^2
+
\left(\frac{2\pi}{L_z}\right)^2
\right]\Phi_{\text{exact}}.
\]
Therefore the matching charge density is
\[
\rho(x,y,z)
=
\varepsilon_0
\left[
\left(\frac{2\pi}{L_x}\right)^2
+
\left(\frac{2\pi}{L_y}\right)^2
+
\left(\frac{2\pi}{L_z}\right)^2
\right]
\Phi_{\text{exact}}(x,y,z).
\]

\subsection*{Why this benchmark is useful}
This choice is useful for three reasons.
\begin{itemize}
\item The fields are smooth, periodic, and analytic.
\item The exact answer is known everywhere in the domain, so error norms are meaningful.
\item The Fourier solver is operating in the setting where it should perform especially well: a rectangular periodic box.
\end{itemize}

\section{Periodic Compatibility and the Zero Mode}
On a periodic domain, Poisson's equation is not solvable for an arbitrary source. Integrating over the domain volume $V$ gives
\[
\int_V \nabla^2 \Phi\, dV
=
-\frac{1}{\varepsilon_0}\int_V \rho\, dV.
\]
The left-hand side vanishes for a periodic potential, so the source must satisfy
\[
\int_V \rho\, dV = 0.
\]
Equivalently, the spatial mean of $\rho$ must be zero. In Fourier language, the $k=0$ mode of the right-hand side must vanish.

For the manufactured cosine benchmark the exact continuous mean is already zero. In the code, the discrete right-hand side is still mean-corrected before the solve. That is a practical numerical safeguard: the FFT-based periodic Poisson solve should not be asked to invert the constant mode.

\section{Fourier-Space View}
The periodic Poisson solve is easiest to understand in Fourier space. If
\[
\nabla^2 \Phi = -\frac{\rho}{\varepsilon_0},
\]
then for each nonzero wavevector $\mathbf k$,
\[
-|\mathbf k|^2 \hat{\Phi}(\mathbf k)
=
-\frac{\hat{\rho}(\mathbf k)}{\varepsilon_0},
\]
so
\[
\hat{\Phi}(\mathbf k)
=
\frac{\hat{\rho}(\mathbf k)}{\varepsilon_0 |\mathbf k|^2},
\qquad
\mathbf k \neq 0.
\]
The constant mode is not inverted. Physically, that reflects the fact that a periodic Poisson problem determines the potential only up to an additive constant and requires zero net charge in the box.

\section{Discrete Interpretation in This Repository}
There are two closely related solvers in the repository.
\begin{itemize}
\item The existing standalone solver in \texttt{codes/cpp/fft\_poisson/fft\_poisson\_3d.cpp}.
\item The AMReX mini-app in \texttt{codes/amrex/fft\_poisson\_benchmark/main.cpp}.
\end{itemize}

They solve the same analytic benchmark, but they do not have to use the same sample locations internally. This matters.

\subsection*{Why the notebook cares about sample coordinates}
The AMReX mini-app writes the coordinate arrays \texttt{x.npy}, \texttt{y.npy}, and \texttt{z.npy} together with the field arrays. The notebook then evaluates the exact manufactured solution on those same sample locations before computing errors.

That avoids a common comparison mistake: subtracting values that come from different discretization conventions such as nodal samples versus cell-centered samples. If the sample locations differ, a direct subtraction can make a correct solver look wrong.

\section{AMReX Data Structures and Their Physical Meaning}
The main AMReX objects in the mini-app correspond directly to parts of the mathematical problem.

\subsection*{\texttt{Geometry}}
The \texttt{Geometry} object defines the rectangular physical domain, the domain extents $(L_x,L_y,L_z)$, and the fact that all three directions are periodic. This is the mathematical box on which the benchmark is posed.

\subsection*{\texttt{BoxArray}}
The \texttt{BoxArray} defines how the logical grid is tiled into rectangular boxes. For this first-pass benchmark, it is kept simple so the comparison/export path stays easy to reason about.

\subsection*{\texttt{DistributionMapping}}
This tells AMReX how the boxes are assigned to ranks. For a clean first correctness comparison, the exporter currently assumes a serial run with a single local box.

\subsection*{\texttt{MultiFab}}
The \texttt{MultiFab} fields hold the discrete arrays:
\begin{itemize}
\item \texttt{rho}: charge density,
\item \texttt{rhs}: the right-hand side of Poisson's equation,
\item \texttt{phi}: the numerical potential,
\item \texttt{phi\_exact}: the exact manufactured potential sampled on the AMReX grid,
\item \texttt{error}: the difference \texttt{phi - phi\_exact}.
\end{itemize}

\subsection*{\texttt{FFT::Poisson}}
This is the AMReX solver wrapper that performs the periodic Poisson solve. Conceptually it is doing the Fourier-space inversion discussed above, but wrapped in AMReX's data structures.

\section{Step-by-Step Computational Workflow}
The mini-app follows a simple sequence.

\subsection*{1. Read input parameters}
The app reads values from \texttt{codes/amrex/fft\_poisson\_benchmark/inputs}. These include the number of cells per direction, the domain lengths, $\varepsilon_0$, the output plotfile name, and the compare-array directory.

\subsection*{2. Build the periodic domain}
The code constructs a cubic index-space domain with \texttt{n\_cell} cells in each direction and marks all directions as periodic.

\subsection*{3. Fill the manufactured fields}
At each AMReX sample point, the code evaluates $\Phi_{\text{exact}}$ and the matching $\rho$. Then it forms the Poisson right-hand side as
\[
\texttt{rhs} = -\rho/\varepsilon_0.
\]

\subsection*{4. Remove the discrete mean of the right-hand side}
Even though the benchmark is analytically mean-free, the code subtracts the discrete mean of the right-hand side. That keeps the FFT solve compatible with periodic Poisson theory at the discrete level.

\subsection*{5. Solve for the numerical potential}
AMReX's \texttt{FFT::Poisson} object computes the discrete periodic solution \texttt{phi}.

\subsection*{6. Compute the pointwise error}
The code forms
\[
\texttt{error} = \texttt{phi} - \texttt{phi\_exact}.
\]
That error field is written out along with the potential and source.

\subsection*{7. Write outputs for two consumers}
The app writes:
\begin{itemize}
\item an AMReX plotfile for visualization and inspection,
\item raw NumPy arrays for notebook-based comparison.
\end{itemize}

The Python helper script
\texttt{scripts/package\_amrex\_compare\_npz.py}
then packages those arrays into
\texttt{amrex\_compare/fft\_poisson\_benchmark.npz},
which the notebook loads automatically if it exists.

\section{How to Build and Run It}
You must run the commands from the repository root:
\begin{verbatim}
/Users/dajuarez4/Computational-Electrodynamics-
\end{verbatim}
not from a subdirectory such as \texttt{Jackson-problems/}.

The AMReX CMake package must already exist somewhere on your machine. The build system needs access to
\begin{verbatim}
AMReXConfig.cmake
\end{verbatim}
either through the AMReX install prefix or directly through the package directory.

A standard configure-and-build sequence is:
\begin{verbatim}
cmake -S codes/amrex/fft_poisson_benchmark \
  -B codes/amrex/fft_poisson_benchmark/build \
  -DAMReX_ROOT=/absolute/path/to/amrex-install

cmake --build codes/amrex/fft_poisson_benchmark/build -j
\end{verbatim}

If you know the exact package directory, this also works:
\begin{verbatim}
cmake -S codes/amrex/fft_poisson_benchmark \
  -B codes/amrex/fft_poisson_benchmark/build \
  -DAMReX_DIR=/absolute/path/to/amrex-install/lib/cmake/AMReX
\end{verbatim}

Then run:
\begin{verbatim}
./codes/amrex/fft_poisson_benchmark/build/amrex_fft_poisson_benchmark \
  codes/amrex/fft_poisson_benchmark/inputs
\end{verbatim}

Then package the notebook input:
\begin{verbatim}
python3 scripts/package_amrex_compare_npz.py \
  --input-dir amrex_compare/fft_poisson_benchmark_amrex \
  --output amrex_compare/fft_poisson_benchmark.npz
\end{verbatim}

\subsection*{A concrete local install layout}
If AMReX is not installed yet, a clean concrete layout on this machine would be:
\begin{verbatim}
AMReX source:   /Users/dajuarez4/src/amrex
AMReX install:  /Users/dajuarez4/.local/amrex
\end{verbatim}

Then the AMReX library itself can be installed with:
\begin{verbatim}
git clone https://github.com/AMReX-Codes/amrex.git /Users/dajuarez4/src/amrex

cmake -S /Users/dajuarez4/src/amrex \
  -B /Users/dajuarez4/src/amrex/build \
  -DCMAKE_BUILD_TYPE=Release \
  -DCMAKE_INSTALL_PREFIX=/Users/dajuarez4/.local/amrex \
  -DAMReX_MPI=NO \
  -DAMReX_OMP=NO \
  -DAMReX_FORTRAN=NO \
  -DAMReX_FFT=YES

cmake --build /Users/dajuarez4/src/amrex/build -j
cmake --install /Users/dajuarez4/src/amrex/build
\end{verbatim}

After that, the benchmark in this repository should be configured with
\begin{verbatim}
-DAMReX_ROOT=/Users/dajuarez4/.local/amrex
\end{verbatim}
or
\begin{verbatim}
-DAMReX_DIR=/Users/dajuarez4/.local/amrex/lib/cmake/AMReX
\end{verbatim}

For CPU FFT builds, note that AMReX uses FFTW. If the AMReX configuration step complains about FFTW, that dependency still needs to be installed before retrying the AMReX build.

\section{Interpreting the Errors You Hit}
Your terminal output showed three different issues, and each one has a different cause.

\subsection*{1. Wrong working directory}
When CMake looked for
\begin{verbatim}
/Users/dajuarez4/Computational-Electrodynamics-/Jackson-problems/codes/amrex/fft_poisson_benchmark
\end{verbatim}
it meant the command was launched from inside \texttt{Jackson-problems/}. The relative path \texttt{codes/amrex/fft\_poisson\_benchmark} only works from the repository root.

\subsection*{2. AMReX is not installed or not discoverable}
The error
\begin{verbatim}
Could not find a package configuration file provided by "AMReX"
\end{verbatim}
means CMake could not locate \texttt{AMReXConfig.cmake}. The most common reasons are:
\begin{itemize}
\item AMReX is not installed yet.
\item The path passed through \texttt{AMReX\_ROOT} is still the literal placeholder \texttt{/path/to/amrex/install}.
\item The package exists, but CMake was not pointed at the correct install prefix or package directory.
\end{itemize}

\subsection*{3. Packaging was attempted before the solver produced arrays}
The Python error about missing
\begin{verbatim}
amrex_compare/fft_poisson_benchmark_amrex/phi.npy
\end{verbatim}
is downstream. It does not indicate a packaging bug. It indicates that the AMReX executable never ran successfully, so there was nothing to package.

\section{What to Check on Your Machine}
Before rebuilding, verify that an AMReX package file actually exists somewhere. On macOS or Linux, one practical search is:
\begin{verbatim}
find /some/search/root -name AMReXConfig.cmake 2>/dev/null
\end{verbatim}
If you find something like
\begin{verbatim}
/absolute/path/to/amrex-install/lib/cmake/AMReX/AMReXConfig.cmake
\end{verbatim}
then either of the following should be valid:
\begin{verbatim}
-DAMReX_ROOT=/absolute/path/to/amrex-install
\end{verbatim}
or
\begin{verbatim}
-DAMReX_DIR=/absolute/path/to/amrex-install/lib/cmake/AMReX
\end{verbatim}

\section{Why This Comparison Is Worth Doing}
The AMReX version is useful even though the current benchmark is simple.
\begin{itemize}
\item It establishes a correct periodic Poisson workflow inside AMReX.
\item It gives a direct apples-to-apples comparison against the existing FFT solver.
\item It creates a foundation for harder problems later: multi-level grids, different source terms, field diagnostics, and eventually problems where AMReX's mesh and parallel infrastructure matter more than a bare FFT code.
\end{itemize}

\section*{Summary}
The benchmark solves a periodic Poisson problem with a manufactured analytic answer, removes the incompatible zero mode, computes the numerical potential with AMReX's FFT-based Poisson solver, and exports data for notebook comparison. Your current blocker is not the physics and not the notebook. The blocker is simply that CMake cannot yet see a real AMReX installation on your machine.

\end{document}
